

\subsection{Proof of Lemma~\ref{lem:bin-dp}} \label{apx:lem:bin-dp}

Let $c= p/(1-p)$. 
We need to find $a < np$ and $b > np$ satisfying the following condition.
\begin{align*}
    \frac{\Pr_{\Bin}[a]}{\Pr_{\Bin+s}[a]} = 
    \frac{ \binom{n}{a} p^a (1-p)^{n-a}} 
         {  \binom{n}{a-s} p^{a-s} (1-p)^{n-a+s}} 
   < \left ( \frac{ n-a} {a -s} \right)^s  \cdot c^s 
   \le e^\ee.
\end{align*}
\begin{align*}
    \frac{\Pr_{\Bin}[b]}{\Pr_{\Bin+s}[b]} = 
    \frac{ \binom{n}{b} p^{b} (1-p)^{n-b}} 
         {  \binom{n}{b-s} p^{b-s} (1-p)^{n-b+s}}
   > \left ( \frac{ n-b} {b} \right)^s  \cdot c^s 
   \ge e^{-\ee}.
\end{align*}

\paragraph{Case 1: $s \ge \ee$.}
We set $a = \frac{np + s(1-p)\cdot e^{\ee/s}}{e^{\ee/s}\cdot (1-p) + p}$ and
$b = \frac{e^{\ee/s} \cdot np}{(1-p) + e^{\ee/s} \cdot p}$. Note these values satisfy the above inequalities. 

To show the second requirement of the tail bound, it suffices to show that $\Pr_{\Bin}[X
\le a+s] = \negl(\kappa)$; the other case holds
similarly.

Let $\mu := np \in \Theta(\kappa)$ and let $d = 1- (a+s)/\mu$. 
By applying the Chernoff bound, we have 
\[ \Pr_{X \from \Bin}[X \le a+s] = \Pr[ X \le (1-d) \mu ] \le
\exp(-d^2\mu/2). \]
%
We will show that we have $d = \Omega(\frac{1}{\lg\lg \kappa})$, which implies that with $\mu \in \Theta(\kappa)$, the above probability is negligible in $\kappa$. 
Let $t = s(1-p)\cdot e^{\ee/s}$ and $u  = e^{\ee/s}\cdot (1-p) + p$. Then, we have $a = \frac{\mu + t}{u}$.
Note that we have $\ee/s \cdot (1-p) + 1 \le u \le e$. We have the following: 
\begin{align*}
    d = \frac{\mu - a - s}{\mu} 
      = \left(1 - \frac{1}{u}\right) - \frac{t/u + s}{\mu}  
      \ge \frac{\ee/s \cdot (1-p)}{e} - \frac{t/u + s}{\mu}  
      = \Omega\left(\frac{1}{\lg\lg \kappa}\right) - \tilde O(1/\kappa).
\end{align*}



\paragraph{Case 2: $s \le \ee$.} 
Let $a = \frac{c}{c+e} \cdot (n+e\ee) < np$ and $b = \frac{c}{c+e^{-1}}
\cdot n > np$. Observe that $\frac{n-a} {b-s} \cdot c = e$ and
$\frac{n-b}{b} \cdot c = e^{-1}$. Therefore, given $s \le \ee$, the above inequalities hold. The
tail bounds specified as the second condition of the lemma can be shown using
the Chernoff bound since $np-a \in \Theta(np) = \Theta(\secparam)$ and so does $b-np$.

\subsection{Proof of Lemma~\ref{lem:good_event}} \label{apx:proof:good_event}
Let $t = \frac{1}{n_R}$,  Let $\mathsf{low} =
1-\left(\frac{1}{2}+\theta\right)^t$ and $\mathsf{high}
= 1-\left(\frac{1}{2}-\theta\right)^t$. 
Then, we have:
\begin{align*}
    &\Pr_{h}[\good_\theta(h, A, I, n_B)]  \\
    & = \Pr_{h} \left[\left(\min h(I) \geq\mathsf{low}\right) \land \left(\min h(I) = \min h(A)\right)\right] \\
    & ~~~~ - \Pr_{h} \left[\left(\min h(I) \geq \mathsf{high}\right) \land \left(\min h(I) = \min h(A)\right)\right]\\
    & = \Pr_{h} \left[\min h(A) \geq\mathsf{low}\right] \cdot \Pr\left[ \min h(I) = \min h(A) \mid \min h(A) \geq\mathsf{low}\right] \\
    & ~~~~ - \Pr_{h} \left[\min h(A) \geq \mathsf{high}\right] \cdot \Pr\left[ \min h(I) = \min h(A) \mid \min h(A) \geq \mathsf{high}\right] \\
    & \ge \left(\frac{1}{2}+\theta\right)^{t \cdot n_A} \cdot
    \frac{n_I}{n_A} - \left(\frac{1}{2}-\theta\right)^{t \cdot n_A}
    \cdot \frac{n_I}{n_A}
\end{align*}

\subsection{Proof of Lemma~\ref{thm:good_hash}} \label{apx:proof:good_hash} 
We can lowerbound $|K_\theta|$ with $\Bino(k-s, p_\theta)$. Recall $s \in O(\lg\lg\kappa)$. 
Let $\mu = (k-s) p_\theta = \Omega(\secparam).$ Applying the Chernoff
Bound, we have 
%
$\Pr\left[ |K_\theta| \le (1 - 1/3) \mu \right] \le \exp(-\mu/18) \le \negl(\secparam). \qedhere$





\subsection{Proof of Lemma~\ref{lem:pb-dp}} \label{apx:proof:pb-dp}
\paragraph{Hockey stick divergence.}
We first review hockey stick divergence~\cite{SasVer15,Ost02,ICALP:BarOlm13}. 
%
%In this paper, we only consider discrete sample space, although many of
%arguments naturally extend to the continuous space. 
% 
%\begin{definition}
The hockey-stick divergence between two probability measures $P,Q$ over $Z$ is defined as:
\begin{equation*}
    \hsd{\alpha}{X}{Y} = \sup_{S \subseteq Z}(X(S) - \alpha Y(S)) = \sum_{z \in Z} [(X(z) - \alpha Y(z)]_+,
\end{equation*}
where $\alpha \geq 1$ and $[x]_+ = \max\{x,0\}$.
We observe that the following holds directly from the definition of the hockey stick divergence.
\begin{corollary}
For any probability measures  $X,Y$ over $Z$ and for any $\ee, \dd$, it holds 
\[ X \aed Y \mbox{ if and only if }  \hsd{e^\ee}{X}{Y} \leq \dd \mbox{ and } \hsd{e^\ee}{Y}{X} \leq \dd.\]
\end{corollary}

\iffalse
Therefore, the hockey-stick divergence can be used to analyze the privacy loss.
%
\paragraph{The hockey-stick divergence as a $f$-divergence.}
It is known that the hockey-stick divergence is a kind of $f$-divergence
~\cite{ICALP:BarOlm13}.  Therefore, the hockey-stick divergence
satisfies the joint convexity and data processing
inequality~\cite{ICALP:BarOlm13,BBG18}.

\begin{lemma}[Joint convexity] 
\label{lem:convex}
For all $0 \leq \lambda \leq 1$ and $\alpha\geq 1$, any probability
measures $X_1,X_2,Y_1,Y_2$ satisfy 
%
\begin{equation*}
    \hsd{\alpha}{\lambda X_1 + (1-\lambda)X_2}{\lambda Y_1 +
    (1-\lambda)Y_2} \leq \lambda \hsd{\alpha}{X_1}{Y_1} +
    (1-\lambda)\hsd{\alpha}{X_2}{Y_2}.  
\end{equation*}
\end{lemma}

The data processing inequality property guarantees that post-processing
cannot increase the divergence.
\begin{lemma}[Data processing inequality.] \label{lem:dpe}
For any probability measures $X, Y$ over $Z$ and any function $g$ with domain $Z$, we have 
\begin{equation*}
    \hsd{\alpha}{g(X)}{g(Y)} \le \hsd{\alpha}{X}{Y}.
\end{equation*}
\end{lemma}
\fi

\paragraph{Proof of Lemma 6.}
Let $\etan = 1/2 - \theta$ and $\etap = 1/2 + \theta$. 
For brevity, we let $C$ denote $\PB(n,p_J)$. For any distribution
$\DDD$, let $P_\DDD$ denote the probability measure with respect to
$\DDD$.
We first show that for any $\ee>0$, it holds that  
$\hsd{e^\ee}{P_{C}}{P_{C + 1}}$ is at most 
\begin{align*}
   \max\left(\hsd{e^\ee}{P_{\Bino(\lceil\frac{n}{2}\rceil, \etap)}}{P_{\Bino(\lceil\frac{n}{2}\rceil, \etap)+1}},
    \hsd{e^\ee}{P_{\Bino(\lceil\frac{n}{2}\rceil, \eta_{-\theta})}}{P_{\Bino(\lceil\frac{n}{2}\rceil, \eta_{-\theta})+1}} \right).
\end{align*}

%To show the above, we follow a similar structure to the proof of \cite[Theorem
%3.3]{COK22}, which analyzes the divergence of a Poisson binomial mechanism.
 %We also rely on the joint convexity\footnote{\cite{COK22}
%actually uses
%joint quasi-convexity, which is implied by joint convexity but not vice
%versa.} (Lemma \ref{lem:convex}) and data processing inequality (Lemma
%\ref{lem:dpe}) to upper bound the hockey-stick divergence of two Poisson
%binomial distributions with that of two binomial distributions. 

We start with an upper bound of the hockey-stick divergence is reached at
extreme points. We rely on the results in \cite{COK22}. Although they use the
Renyi divergence, their results are general enough to be applied to any $f$-divergence.

\begin{lemma} (\cite[Lemma 3.5]{COK22})
\begin{align}
     \hsd{e^\ee}{P_{C}}{P_{C + 1}}
     \leq  \max_{j \in [n]} \hsd{e^\ee}{P_{\Bino(j, \eta_{-\theta}) + \Bino(n-j,\eta_{+\theta})}}{P_{\Bino(j, \eta_{-\theta}) + \Bino(n-j,\eta_{+\theta})+1}}.\label{eqa:twobinom}
\end{align}
\end{lemma}

\iffalse
\begin{proof}
Note that there is $\lambda \in [0, 1]$ such that $p_{n} = \lambda\left(\eta_{-\theta}\right) + (1-\lambda)\left(\eta_{+\theta}\right)$. 
Let $C_{-n} = \sum_{j=1}^{n-1} c_j$. 
Denote the convolution $\conv(a, a', b) = a \cdot (1-b) + a' \cdot b.$
Then, we have 
\begin{align*}
\Pr[C = \ell] & = \conv( \Pr[C_{-n} = \ell], \Pr[C_{-n} = \ell-1], p_{n}) \\
 & = \conv\left(\Pr[C_{-n} = \ell], \Pr[C_{-n} = \ell-1], \lambda \cdot
 \eta_{-\theta} + (1-\lambda) \cdot \eta_{+\theta}\right) \\
 & = \lambda \conv(\Pr[C_{-n} = \ell], \Pr[C_{-n} = \ell-1], \eta_{-\theta}) \\
 & ~~~~~ + (1-\lambda) \conv(\Pr[C_{-n} = \ell], \Pr[C_{-n} = \ell-1], \eta_{+\theta}) \\
 & = \lambda \Pr[C_{-n} + \Ber(\eta_{-\theta}) = \ell] + 
    (1-\lambda) \Pr[ C_{-n} + \Ber(\eta_{+\theta}) = \ell] 
\end{align*}
%
In other words, we have 
\[P_C = \lambda P_{C_{-n} + \Ber(\eta_{-\theta})} + (1-\lambda) P_{C_{-n} + \Ber(\eta_{+\theta})}.\]
Now using Lemma~\ref{lem:convex}, we have
\begin{align*}
& \hsd{e^\ee}{P_{C}}{P_{C + 1}} \\
& ~~ \leq\lambda \hsd{e^\ee}{ 
  P_{C_{-n} + \Ber(\eta_{-\theta})}}{ 
  P_{C_{-n} + \Ber(\eta_{-\theta}) + 1}}
  + (1-\lambda) \hsd{e^\ee}
    { P_{C_{-n} + \Ber(\eta_{+\theta})}}
    { P_{C_{-n} + \Ber(\eta_{+\theta}) + 1}}\\
& ~~ \leq \max\left\{ \hsd{e^\ee}
    { P_{C_{-n} + \Ber(\eta_{-\theta})} }
    { P_{C_{-n} + \Ber(\eta_{-\theta}) + 1}},
    \hsd{e^\ee}
    { P_{C_{-n} + \Ber(\eta_{+\theta})} }
    { P_{C_{-n} + \Ber(\eta_{+\theta}) + 1}} \right\}
\end{align*}
Repetitively applying the joint convexity on the remaining $n-1$ random variables yields the following
\begin{align*}
     &\hsd{e^\ee}{P_{Y}}{P_{Y+ 1}} \leq \max_{j\in [n]} \hsd{e^\ee}{P_{N_j}}{P_{N_j+1}},
\end{align*}
where $N_j \sim \Bino(j, \eta_{-\theta}) + \Bino(n-j,\eta_{+\theta})$.
\end{proof}
\fi

Next, we apply data processing inequality to simplify (\ref{eqa:twobinom}) from the above lemma. 

\begin{lemma} (\cite[Lemma 3.6]{COK22})
(\ref{eqa:twobinom}) is upper bounded by
\begin{align}
\max \left(    \hsd{e^\ee}{P_{\Bino(\lceil\frac{n}{2}\rceil, \etap)}}{P_{\Bino(\lceil\frac{n}{2}\rceil, \etap)+1}},
    \hsd{e^\ee}{P_{\Bino(\lceil\frac{n}{2}\rceil, \eta_{-\theta})}}{P_{\Bino(\lceil\frac{n}{2}\rceil, \eta_{-\theta})+1}}\right).
\end{align}
\end{lemma}
\iffalse
\begin{proof}
The proof is similar to \cite[Lemma 3.6]{COK22}.  We apply the data
processing inequality. More specifically, consider two distributions $X,
Y$ and then adding a binomial random variable $Z$ as the post-processing
step. Then, the addition of $Z$ doesn't increase the divergence. In
other words, \[\hsd{e^\ee}{P_X}{P_Y} \ge \hsd{e^\ee}{P_{X+Z}}{P_{Y+Z}}.\] 

Now, consider (\ref{eqa:twobinom}) and let $j^*$ be the index that leads
to the maximum. 
\begin{itemize}
\item If $j^* \le n/2$, we have 
\begin{align*}
& \hsd{e^\ee}{P_{\Bino(\lceil\frac{n}{2}\rceil, \eta_{+\theta})}}
{P_{\Bino(\lceil\frac{n}{2}\rceil, \eta_{+\theta})+1}}  \\
& ~~~\ge \hsd{e^\ee}{P_{\Bino(j^*, \eta_{-\theta}) +
\Bino(n-j^*,\eta_{+\theta})}}{P_{\Bino(j^*, \eta_{-\theta}) +
\Bino(n-j^*,\eta_{+\theta})+1}}.
\end{align*}
This is because $n/2 \le n - j^*$. In this case, $Z = \Bino(j^*, \eta_{-\theta}) +
\Bino(n-j^*-\lceil\frac{n}{2}\rceil,\eta_{+\theta})$.

\item Likewise, if $j^* \ge n/2$, we have 
\begin{align*}
&  \hsd{e^\ee}{P_{\Bino(\lceil\frac{n}{2}\rceil, \eta_{-\theta})}}
{P_{\Bino(\lceil\frac{n}{2}\rceil, \eta_{-\theta})+1}} \\
& ~~~ \ge \hsd{e^\ee}{P_{\Bino(j^*, \eta_{-\theta}) +
\Bino(n-j^*,\eta_{+\theta})}}{P_{\Bino(j^*, \eta_{-\theta}) +
\Bino(n-j^*,\eta_{+\theta})+1}}.
\end{align*}
\end{itemize} 
\end{proof}
\fi


We extend the above to upper bound the hockey-stick divergence between probability
measures differed by an integer amount greater than 1, i.e.,  $P_{C}$
and $P_{C+s}$ for $s>1$. 

\begin{corollary}\label{col:APBM}
For any $\ee>0$, it holds that 
    $\hsd{e^\ee}{P_{C}}{P_{C + s}}$ is at most 
\begin{align*}
    \max \left( 
    \hsd{e^\ee}{P_{\Bino(\lceil\frac{n}{2}\rceil, \etap)}}{P_{\Bino(\lceil\frac{n}{2}\rceil, \etap)+s}},
    \hsd{e^\ee}{P_{\Bino(\lceil\frac{n}{2}\rceil, \eta_{-\theta})}}{P_{\Bino(\lceil\frac{n}{2}\rceil, \eta_{-\theta})+s}} \right).
\end{align*}
\end{corollary}

Finally, to give a bound on the divergence, we can apply Lemma~\ref{lem:bin-dp} to argue that the binomial distribution hides the small sensitivity. Specifically, as $\lceil\frac{n}{2}\rceil \in \Theta(\kappa)$ and $s =\lg\lg\kappa$, we can claim $(\ee, \dd)$-DDP with $\dd = \negl(\kappa)$.

Similarly, it holds that $\hsd{e^\ee}{P_{C+s}}{P_{C}} \le
\negl(\kappa)$. \qed 
